\section{Continuous maps and abelian sheaves}
\label{section-abelian-presheaves-functorial}

\noindent
Let $f : X \to Y$ be a continuous map.
We claim there are functors
\begin{eqnarray*}
f_* : \textit{PAb}(X) & \longrightarrow & \textit{PAb}(Y) \\
f_* : \textit{Ab}(X) & \longrightarrow & \textit{Ab}(Y) \\
f_p : \textit{PAb}(Y) & \longrightarrow & \textit{PAb}(X) \\
f^{-1} : \textit{Ab}(Y) & \longrightarrow & \textit{Ab}(X)
\end{eqnarray*}
with similar properties to their counterparts in
Section \ref{section-presheaves-functorial}.
To see this we argue in the following way.

\medskip\noindent
Each of the functors will be constructed in the same
way as the corresponding functor in
Section \ref{section-presheaves-functorial}.
This works because all the colimits in that section
are directed colimits (but we will work through it below).

\medskip\noindent
First off, given an abelian presheaf $\mathcal{F}$ on $X$ and
an abelian presheaf $\mathcal{G}$ on $Y$ we define
\begin{eqnarray*}
f_*\mathcal{F}(V) & = & \mathcal{F}(f^{-1}(V)) \\
f_p\mathcal{G}(U) & = & \colim_{f(U) \subset V} \mathcal{G}(V)
\end{eqnarray*}
as abelian groups. The restriction mappings are the same as
the restriction mappings for presheaves of sets (and they are
all homomorphisms of abelian groups).

\medskip\noindent
The assignments $\mathcal{F} \mapsto f_*\mathcal{F}$ and
$\mathcal{G} \to f_p\mathcal{G}$ are functors on
the categories of presheaves of abelian groups.
This is clear, as (for example) a map of abelian presheaves
$\mathcal{G}_1 \to \mathcal{G}_2$ gives rise to a map of
directed systems
$\{\mathcal{G}_1(V)\}_{f(U) \subset V} \to
\{\mathcal{G}_2(V)\}_{f(U) \subset V}$
all of whose maps are homomorphisms
and hence gives rise to a homomorphism of abelian groups
$f_p\mathcal{G}_1(U) \to f_p\mathcal{G}_2(U)$.

\medskip\noindent
The functors $f_*$ and $f_p$ are adjoint on the
category of presheaves of abelian groups, i.e., we have
$$
\Mor_{\textit{PAb}(X)}(f_p\mathcal{G}, \mathcal{F})
=
\Mor_{\textit{PAb}(Y)}(\mathcal{G}, f_*\mathcal{F}).
$$
To prove this, note that the map
$i_\mathcal{G} : \mathcal{G} \to f_* f_p\mathcal{G}$ from the proof
of Lemma \ref{lemma-pullback-presheaves}
is a map of abelian presheaves. Hence if
$\psi : f_p\mathcal{G} \to \mathcal{F}$
is a map of abelian presheaves, then the corresponding map
$\mathcal{G} \to f_*\mathcal{F}$ is the map
$f_*\psi \circ i_\mathcal{G} :
\mathcal{G} \to f_* f_p \mathcal{G} \to f_* \mathcal{F}$
is also a map of abelian presheaves. For the other direction
we point out that the map
$c_\mathcal{F} : f_p f_* \mathcal{F} \to \mathcal{F}$
from the proof of Lemma \ref{lemma-pullback-presheaves} is a
map of abelian presheaves as well (since it is made out of restriction
mappings of $\mathcal{F}$ which are all homomorphisms). Hence
given a map of abelian presheaves $\varphi : \mathcal{G} \to f_*\mathcal{F}$
the map
$c_\mathcal{F} \circ f_p\varphi : f_p\mathcal{G} \to \mathcal{F}$
is a map of abelian presheaves as well. Since these constructions
$\psi \mapsto f_*\psi$ and $\varphi \mapsto c_\mathcal{F} \circ f_p\varphi$
are inverse to each other as constructions on maps of presheaves of sets
we see they are also inverse to each other on maps of abelian presheaves.

\medskip\noindent
If $\mathcal{F}$ is an abelian sheaf on $Y$, then $f_*\mathcal{F}$
is an abelian sheaf on $X$. This is true because of the definition
of an abelian sheaf and because this is true for sheaves of sets,
see Lemma \ref{lemma-pushforward-sheaf}. This defines the functor
$f_*$ on the category of abelian sheaves.

\medskip\noindent
We define $f^{-1}\mathcal{G} = (f_p\mathcal{G})^\#$ as before.
Adjointness of $f_*$ and $f^{-1}$ follows formally as in
the case of presheaves of sets. Here is the argument:
\begin{eqnarray*}
\Mor_{\textit{Ab}(X)}(f^{-1}\mathcal{G}, \mathcal{F})
& = &
\Mor_{\textit{PAb}(X)}(f_p\mathcal{G}, \mathcal{F}) \\
& = &
\Mor_{\textit{PAb}(Y)}(\mathcal{G}, f_*\mathcal{F}) \\
& = &
\Mor_{\textit{Ab}(Y)}(\mathcal{G}, f_*\mathcal{F})
\end{eqnarray*}

\begin{lemma}
\label{lemma-pullback-abelian-stalk}
Let $f : X \to Y$ be a continuous map.
\begin{enumerate}
\item Let $\mathcal{G}$ be an abelian presheaf on $Y$.
Let $x \in X$. The bijection
$\mathcal{G}_{f(x)} \to (f_p\mathcal{G})_x$ of
Lemma \ref{lemma-stalk-pullback-presheaf} is an isomorphism of abelian groups.
\item Let $\mathcal{G}$ be an abelian sheaf on $Y$.
Let $x \in X$. The bijection
$\mathcal{G}_{f(x)} \to (f^{-1}\mathcal{G})_x$ of
Lemma \ref{lemma-stalk-pullback} is an isomorphism of abelian groups.
\end{enumerate}
\end{lemma}

\begin{proof}
Omitted.
\end{proof}

\noindent
Given a continuous map $f : X \to Y$ and sheaves of abelian
groups $\mathcal{F}$ on $X$, $\mathcal{G}$ on $Y$, the notion
of an {\it $f$-map $\mathcal{G} \to \mathcal{F}$
of sheaves of abelian groups} makes sense. We can just define
it exactly as in Definition \ref{definition-f-map} (replacing maps
of sets with homomorphisms of abelian groups) or we can
simply say that it is the same as a map of abelian
sheaves $\mathcal{G} \to f_*\mathcal{F}$. We will use this
notion freely in the following. The group of $f$-maps between
$\mathcal{G}$ and $\mathcal{F}$ will be in canonical bijection
with the groups
$\Mor_{\textit{Ab}(X)}(f^{-1}\mathcal{G}, \mathcal{F})$
and
$\Mor_{\textit{Ab}(Y)}(\mathcal{G}, f_*\mathcal{F})$.

\medskip\noindent
Composition of $f$-maps is defined in exactly the
same manner as in the case of $f$-maps of sheaves of
sets. In addition, given an $f$-map $\mathcal{G} \to \mathcal{F}$
as above, the induced maps on stalks
$$
\varphi_x : \mathcal{G}_{f(x)} \longrightarrow \mathcal{F}_x
$$
are abelian group homomorphisms.
